\section{Incidence matrices}

\subsection{Vectors and matrices}

Linear algebra encodes Petri nets: the net becomes an integer matrix, markings become vectors, firing becomes vector addition, and a firing sequence becomes an integer vector counting its transitions. The invariants and soundness conditions of later chapters build on this encoding.\footnote{Desel J., Esparza J., \emph{Free Choice Petri Nets}, Cambridge University Press, 1995, \url{https://www7.in.tum.de/~esparza/bookfc.html}.}

\begin{definitionbox}{Vectors, products, linear systems}
For a finite set $E = \{e_1, \ldots, e_n\}$, any map $v : E \to \mathbb{Q}$ (or $\mathbb{N}$, $\mathbb{Z}$) is a vector indexed by $E$; we use the same symbol for the row and the column form, the orientation being clear from the multiplication performed. The \emph{scalar product} of two equal-length vectors is $\mathbf{x} \cdot \mathbf{y} = \sum_i x_i y_i$ (row $\times$ column); a matrix-by-vector product computes one scalar product per row, a vector-by-matrix product one per column. A \emph{homogeneous linear system} is a matrix equation $A \cdot \mathbf{x} = \mathbf{0}$; its non-zero solutions are the invariants of the next chapter.
\end{definitionbox}

\begin{definitionbox}{Markings as vectors; orderings}
A marking $M : P \to \mathbb{N}$ \emph{is} a vector indexed by the places: $M = [M(p_1)\; \cdots\; M(p_n)]$. For vectors over $E$:
\[
  \mathbf{v} \leq \mathbf{w} \iff \forall e.\ v(e) \leq w(e); \qquad
  \mathbf{v} < \mathbf{w} \iff \mathbf{v} \leq \mathbf{w} \wedge \exists e.\ v(e) < w(e); \qquad
  \mathbf{v} \prec \mathbf{w} \iff \forall e.\ v(e) < w(e).
\]
$\mathbf{0}$ is the all-zeros vector of the appropriate length. The two strict orders differ: $<$ allows ties, $\prec$ forbids them: e.g., $[1\,2\,0] < [2\,3\,0]$ holds, but $[0\,0\,0] \prec [2\,0\,1]$ fails on the tied component, while $[0\,0\,0] \prec [2\,3\,1]$ holds. Neither $[3\,2\,0] \leq [2\,3\,1]$ nor $[3\,2\,0] \geq [2\,3\,1]$ holds: vector orders are partial.
\end{definitionbox}

\subsection{The incidence matrix}

Can net computations be expressed by matrix multiplication? Yes, with a caveat: the \emph{change} of marking caused by firing $t$ is structural (it depends only on the arcs of $t$, not on the marking at which $t$ fires) so a single integer per (place, transition) pair captures it; but the \emph{enabling} condition is not linear and will be lost. Four connection cases exhaust the possibilities:
\begin{center}
\begin{tikzpicture}[node distance=0.5cm and 0.9cm]
  \node[bpmn event, label=above:{\scriptsize $p$}] (p1) {};
  \node[bpmn task, rounded corners=0pt, minimum width=0.4cm, minimum height=0.4cm, right=of p1, label=above:{\scriptsize $t$}] (t1) {};
  \node[font=\scriptsize, below=2pt of t1, xshift=-0.55cm] {no arc: $0$};
  \begin{scope}[xshift=3.2cm]
  \node[bpmn event, label=above:{\scriptsize $p$}] (p2) {};
  \node[bpmn task, rounded corners=0pt, minimum width=0.4cm, minimum height=0.4cm, right=of p2, label=above:{\scriptsize $t$}] (t2) {};
  \draw[bpmn flow] (p2) -- (t2);
  \node[font=\scriptsize, below=2pt of t2, xshift=-0.55cm] {input: $-1$};
  \end{scope}
  \begin{scope}[xshift=6.4cm]
  \node[bpmn event, label=above:{\scriptsize $p$}] (p3) {};
  \node[bpmn task, rounded corners=0pt, minimum width=0.4cm, minimum height=0.4cm, right=of p3, label=above:{\scriptsize $t$}] (t3) {};
  \draw[bpmn flow] (t3) -- (p3);
  \node[font=\scriptsize, below=2pt of t3, xshift=-0.55cm] {output: $+1$};
  \end{scope}
  \begin{scope}[xshift=9.6cm]
  \node[bpmn event, label=above:{\scriptsize $p$}] (p4) {};
  \node[bpmn task, rounded corners=0pt, minimum width=0.4cm, minimum height=0.4cm, right=of p4, label=above:{\scriptsize $t$}] (t4) {};
  \draw[bpmn flow] (p4) to[bend right=25] (t4);
  \draw[bpmn flow] (t4) to[bend right=25] (p4);
  \node[font=\scriptsize, below=2pt of t4, xshift=-0.55cm] {self-loop: $0$};
  \end{scope}
\end{tikzpicture}
\end{center}

\begin{definitionbox}{Incidence matrix}
For a net $N = (P, T, F)$ the \textbf{incidence matrix} $\mathbf{N} : (P \times T) \to \{-1, 0, +1\}$ is
\[
\mathbf{N}(p, t) = \begin{cases}
 -1 & \text{if } (p, t) \in F \wedge (t, p) \notin F, \\
 +1 & \text{if } (p, t) \notin F \wedge (t, p) \in F, \\
  0 & \text{otherwise (no arc, or self-loop).}
\end{cases}
\]
With $|P| = n$, $|T| = m$, $\mathbf{N}$ is an $n \times m$ integer matrix: sparse, since each transition touches few places. The third case lumps ``no connection'' and ``self-loop'' together: same \emph{net effect}, different enabling (a self-loop requires $M(p) \geq 1$). That lost information is the limit of structural analysis. The column $\mathbf{t}_j = \mathbf{N}(\cdot, t_j)$ is the effect-vector of $t_j$ on every place; the row $\mathbf{p}_i = \mathbf{N}(p_i, \cdot)$ is the input/output profile of $p_i$ across all transitions.
\end{definitionbox}

\begin{examplebox}{Running example: the vending machine}
Five places ($p_1$ \emph{candy storage}, $p_2$ \emph{request for refill}, $p_3$ \emph{ready for coin}, $p_4$ \emph{holding}, $p_5$ \emph{ready to dispense}) and five transitions: $t_1$ \emph{refill} ($p_2 \to p_1$), $t_2$ \emph{dispense candy} ($p_1 + p_5 \to p_2 + p_3$), $t_3$ \emph{insert coin} ($p_3 \to p_4$), $t_4$ \emph{accept coin} ($p_4 \to p_5$), $t_5$ \emph{reject coin} ($p_4 \to p_3$). Initial marking $M_0 = 4p_1 + p_3$.
\begin{center}
\begin{tikzpicture}
  \node[bpmn event, label=above:{\scriptsize $p_1$ candy storage}] (p1) at (0,1.6) {};
  \fill[accent] ([xshift=-2.5pt,yshift=2.5pt]p1.center) circle (1.3pt);
  \fill[accent] ([xshift=2.5pt,yshift=2.5pt]p1.center) circle (1.3pt);
  \fill[accent] ([xshift=-2.5pt,yshift=-2.5pt]p1.center) circle (1.3pt);
  \fill[accent] ([xshift=2.5pt,yshift=-2.5pt]p1.center) circle (1.3pt);
  \node[bpmn event, label=below:{\scriptsize $p_2$ request for refill}] (p2) at (0,0) {};
  \node[bpmn task, rounded corners=0pt, font=\tiny\sffamily] (t1) at (-1.8,0.8) {refill};
  \node[bpmn task, rounded corners=0pt, font=\tiny\sffamily, align=center] (t2) at (1.8,0.8) {dispense\\ candy};
  \node[bpmn event, label=above:{\scriptsize $p_3$ ready for coin}] (p3) at (3.8,1.6) {};
  \fill[accent] (p3.center) circle (1.5pt);
  \node[bpmn task, rounded corners=0pt, font=\tiny\sffamily, align=center] (t3) at (5.6,1.6) {insert\\ coin};
  \node[bpmn event, label=right:{\scriptsize $p_4$ holding}] (p4) at (7.0,0.8) {};
  \node[bpmn task, rounded corners=0pt, font=\tiny\sffamily, align=center] (t4) at (5.6,0) {accept\\ coin};
  \node[bpmn event, label=below:{\scriptsize $p_5$ ready to dispense}] (p5) at (3.8,0) {};
  \node[bpmn task, rounded corners=0pt, font=\tiny\sffamily, align=center] (t5) at (5.6,2.9) {reject\\ coin};
  \draw[bpmn flow] (p2) to[bend left=20] (t1);
  \draw[bpmn flow] (t1) to[bend left=20] (p1);
  \draw[bpmn flow] (p1) to[bend left=15] (t2);
  \draw[bpmn flow] (t2) to[bend left=15] (p2);
  \draw[bpmn flow] (t2) -- (p3);
  \draw[bpmn flow] (p3) -- (t3);
  \draw[bpmn flow] (t3) -- (p4);
  \draw[bpmn flow] (p4) -- (t4);
  \draw[bpmn flow] (t4) -- (p5);
  \draw[bpmn flow] (p5) -- (t2);
  \draw[bpmn flow] (p4) -- (t5);
  \draw[bpmn flow] (t5) -- (p3);
\end{tikzpicture}
\end{center}
Filling the matrix column by column from the arc structure (e.g., \emph{refill} consumes a refill request and produces a unit of stock):
\[
\mathbf{N} = \begin{bmatrix}
+1 & -1 & 0 & 0 & 0 \\
-1 & +1 & 0 & 0 & 0 \\
 0 & +1 & -1 & 0 & +1 \\
 0 &  0 & +1 & -1 & -1 \\
 0 & -1 & 0 & +1 & 0
\end{bmatrix}
\qquad
\begin{array}{l}
\text{rows: } p_1, \ldots, p_5 \\
\text{columns: } t_1 \text{ refill}, t_2 \text{ dispense}, t_3 \text{ insert}, \\
\qquad t_4 \text{ accept}, t_5 \text{ reject}
\end{array}
\]
Reading the matrix: columns $t_4$ and $t_5$ both carry $-1$ on $p_4$, so accept and reject compete for the same holding token, a structural XOR. Every column sums to zero: each transition moves tokens without creating or destroying them, a property the next chapter names \emph{conservativeness}.
\end{examplebox}

\begin{definitionbox}{Firing in vector notation}
Firing $t_j$ at $M$ produces $M' = M + \mathbf{t}_j$, the $j$-th column added to the marking vector. On the vending machine, the sequence $M_0 \xrightarrow{t_3} M_1 \xrightarrow{t_4} M_2 \xrightarrow{t_2} M_3$ unrolls as
\[
\begin{bmatrix} 4 \\ 0 \\ 1 \\ 0 \\ 0 \end{bmatrix}
\xrightarrow{+\,\mathbf{t}_3}
\begin{bmatrix} 4 \\ 0 \\ 0 \\ 1 \\ 0 \end{bmatrix}
\xrightarrow{+\,\mathbf{t}_4}
\begin{bmatrix} 4 \\ 0 \\ 0 \\ 0 \\ 1 \end{bmatrix}
\xrightarrow{+\,\mathbf{t}_2}
\begin{bmatrix} 3 \\ 1 \\ 1 \\ 0 \\ 0 \end{bmatrix}:
\]
one candy sold, a refill request raised, the machine back at \emph{ready for coin}. The final marking is $M_0 + \mathbf{t}_3 + \mathbf{t}_4 + \mathbf{t}_2$: the order of the summands is irrelevant, only their multiset matters. That multiset is the Parikh vector.
\end{definitionbox}

A scalar-product warm-up that will matter later: dotting the row of $p_3$ with the column $\mathbf{y} = [1, 1, 2, 0, 1]^T$,
\[
\begin{bmatrix} 0 & 1 & -1 & 0 & 1 \end{bmatrix} \cdot \begin{bmatrix} 1 \\ 1 \\ 2 \\ 0 \\ 1 \end{bmatrix} = 0 + 1 - 2 + 0 + 1 = 0,
\]
and computing the full product $\mathbf{N} \cdot \mathbf{y} = [0, 0, 0, 1, -1]^T$: the first three rows cancel, the last two are opposite. Vectors $\mathbf{y}$ that annul such products are the invariants of the next chapter.

\subsection{The marking equation}

\begin{definitionbox}{Parikh vector}
For $\sigma \in T^\star$, the \textbf{Parikh vector} $\vec{\sigma} : T \to \mathbb{N}$ maps each transition to its number of occurrences in $\sigma$. Inductively: $\vec{\varepsilon} = \mathbf{0}$; $\vec{t} = \mathbf{e}_t$ (the standard basis vector with a single $1$ at position $t$); $\vec{\sigma t} = \vec{\sigma} + \vec{t}$. Hence $\vec{\sigma_1 \sigma_2} = \vec{\sigma_1} + \vec{\sigma_2}$: the Parikh map forgets order and retains multiplicity: permutations share the same vector.
\end{definitionbox}

On the vending machine, $M_0 \xrightarrow{\sigma} 3p_1 + p_2 + p_3$ with $\sigma = t_3 t_5 t_3 t_4 t_2$ (insert, reject, re-insert, accept, dispense) has $\vec{\sigma} = [0, 1, 2, 1, 1]$; the longer run $M_0 \xrightarrow{\sigma'} 2p_1 + 2p_2 + p_4$ with $\sigma' = t_3 t_4 t_2\, t_3 t_4 t_2\, t_3 t_5 t_3$ has $\vec{\sigma'} = [0, 2, 4, 2, 1]$: twice the dispense events, the same single rejection.

Since $\vec{t}_j = \mathbf{e}_j$ and multiplying a matrix by $\mathbf{e}_j$ selects its $j$-th column, $\mathbf{N} \cdot \vec{t}_j = \mathbf{t}_j$, so the single-step rule reads $M \xrightarrow{t} M' \Rightarrow M' = M + \mathbf{N} \cdot \vec{t}$. The generalisation to sequences is the central lemma of the chapter.

\begin{lemma}[Marking equation]
If $M \xrightarrow{\sigma} M'$ then $M' = M + \mathbf{N} \cdot \vec{\sigma}$.
\end{lemma}

\textit{Proof.} Induction on $|\sigma|$. Base $\sigma = \varepsilon$: $M' = M$ and $\mathbf{N} \cdot \mathbf{0} = \mathbf{0}$. Step $\sigma = \sigma' t$ with $M \xrightarrow{\sigma'} M'' \xrightarrow{t} M'$:
\[
M' = M'' + \mathbf{N} \cdot \vec{t} = M + \mathbf{N} \cdot \vec{\sigma'} + \mathbf{N} \cdot \vec{t} = M + \mathbf{N} \cdot (\vec{\sigma'} + \vec{t}) = M + \mathbf{N} \cdot \vec{\sigma' t}. \qquad \square
\]

On the vending machine, $M_0 + \mathbf{N} \cdot [0,1,2,1,1]^T = [4,0,1,0,0]^T + [-1,+1,0,0,0]^T = [3,1,1,0,0]^T$: the marking $3p_1 + p_2 + p_3$ recovered without simulating. The implication is one-way: it computes a sequence's effect but says nothing about whether the sequence was enabled.

\subsection{Monotonicity and its consequences}

The monotonicity proof below and its corollary reach into infinite firing sequences, whose enabledness reduces to the finite case through the prefix characterisation proved in the automata-to-nets chapter: $M \xrightarrow{\sigma}$ holds for an infinite $\sigma$ iff every finite prefix of $\sigma$ is enabled at $M$.

\begin{lemma}[Monotonicity]
If $M \xrightarrow{\sigma} M'$ then $M + L \xrightarrow{\sigma} M' + L$ for every marking $L \geq \mathbf{0}$; the same holds for infinite $\sigma$ (if $M \xrightarrow{\sigma}$ then $M + L \xrightarrow{\sigma}$).
\end{lemma}

\textit{Proof.} Finite case by induction on $|\sigma|$: the empty sequence is enabled anywhere; for $\sigma = \sigma' t$, the induction gives $M + L \xrightarrow{\sigma'} M'' + L$, where $t$ is still enabled (extra tokens only help) and fires to $M'' + L + \mathbf{N} \cdot \vec{t} = M' + L$. Infinite case: by the prefix characterisation it suffices that every finite prefix be enabled at $M + L$, which the finite case provides. \hfill$\square$

Two readings: what can be done with fewer resources can be done with more; and the surplus survives untouched. On the vending machine, $M_0 \xrightarrow{\sigma} [3,1,1,0,0]^T$ with $\sigma = t_3 t_5 t_3 t_4 t_2$ lifts to $M_0 + L \xrightarrow{\sigma} [4,2,1,0,0]^T$ for $L = [1,1,0,0,0]^T$: the two extra tokens are spectators.

\begin{corollary}[Repeatable prefix]
If $M \xrightarrow{\sigma} M'$ with $M \le M'$, then $M \xrightarrow{\sigma\sigma\sigma\cdots}$.
\end{corollary}

\textit{Proof.} Every prefix of $\sigma^\omega$ has the form $\sigma^n \sigma'$ with $\sigma'$ a prefix of $\sigma$; induction on $n$ with $L = M' - M \geq \mathbf{0}$: the monotonicity lemma turns $M \xrightarrow{\sigma^n \sigma'}$ into $M' = M + L \xrightarrow{\sigma^n \sigma'}$, and prefixing one $\sigma$ gives $M \xrightarrow{\sigma^{n+1} \sigma'}$. \hfill$\square$

\begin{lemma}[Boundedness]
In a bounded system, $M \in [M_0\rangle$ with $M \ge M_0$ forces $M = M_0$.
\end{lemma}

\textit{Proof.} Let $M = M_0 + L$ and $M_0 \xrightarrow{\sigma} M$. Iterating the monotonicity lemma, $M_0 \xrightarrow{\sigma} M_0 + L \xrightarrow{\sigma} M_0 + 2L \xrightarrow{\sigma} \cdots$, so every $M_0 + kL$ is reachable; boundedness forces this set to be finite, hence $L = \mathbf{0}$. \hfill$\square$

Contrapositive diagnostic: a reachable $M > M_0$ proves unboundedness.

\begin{lemma}[Repetition]
If $M \xrightarrow{\sigma} M'$ and $\sigma$ can be fired arbitrarily many times from $M$, then $M \le M'$.
\end{lemma}

\textit{Proof.} Otherwise some place has $M'(p) = M(p) - k$ with $k > 0$, i.e., $(\mathbf{N} \cdot \vec{\sigma})(p) = -k$; firing $\sigma$ $n = M(p) + 1$ times would give $M(p) - nk < 0$, absurd. \hfill$\square$

A repeatable sequence satisfies $\mathbf{N} \cdot \vec{\sigma} \geq \mathbf{0}$: it is a \emph{cycle of resources}, producing on every place at least what it consumes. The Parikh vectors with $\mathbf{N} \cdot \vec{\sigma} = \mathbf{0}$ form the kernel of $\mathbf{N}$: the \emph{T-invariants} of the next chapter.

\begin{examplebox}{Exercise: diagnosing sequences by matrix algebra}
With $M_0 = 4p_1 + p_3$:
\par\smallskip
\emph{(a)} $\sigma = t_3 t_4 t_2\, t_3 t_5\, t_3 t_4\, t_1 t_2\, t_1\, t_3 t_5\, t_2$ has Parikh vector $\vec{\sigma} = [2, 3, 4, 2, 2]$. The marking equation gives
\[
M_0 + \mathbf{N} \cdot \vec{\sigma} = \begin{bmatrix} 3 \\ 1 \\ 2 \\ 0 \\ \boxed{-1} \end{bmatrix},
\]
a negative entry on $p_5$: impossible. Hence \emph{no} firing sequence with this Parikh vector is enabled at $M_0$: non-fireability proved without simulating a single prefix.
\par\smallskip
\emph{(b)} $\sigma' = t_3 t_4 t_2 t_1$ (insert, accept, dispense, refill) has $\vec{\sigma'} = [1, 1, 1, 1, 0]$ and
\[
\mathbf{N} \cdot \vec{\sigma'} = \mathbf{0}, \qquad \text{so } M_0 + \mathbf{N} \cdot \vec{\sigma'} = M_0.
\]
Direct simulation shows $\sigma'$ is fireable at $M_0$ (each step enabled in turn), and it returns exactly to $M_0 = M_0$: by the repeatable-prefix corollary, $M_0 \xrightarrow{\sigma' \sigma' \cdots}$: the insert/accept/dispense/refill cycle runs forever. Its Parikh vector lies in the kernel of $\mathbf{N}$: a T-invariant, the object of the next chapter.
\end{examplebox}
